\section{Results and Performance Evaluation}
\label{sec:results}

\subsection{Quantitative Benchmark on 1,611 Held-Out Test Patches}
We evaluate reconstruction performance using five standard quantitative remote sensing metrics:
\begin{enumerate}
    \item \textbf{Peak Signal-to-Noise Ratio (PSNR, dB):} Measures global pixel reconstruction fidelity (higher is better).
    \item \textbf{Structural Similarity Index (SSIM) \cite{wang2004image}:} Evaluates structural texture, luminance, and contrast preservation (higher is better, max 1.0).
    \item \textbf{Spectral Angle Mapper (SAM, degrees):} Computes the average angular difference between predicted and reference spectral vectors, assessing color/radiometric fidelity (lower is better, 0.0$^\circ$ is perfect).
    \item \textbf{Relative Dimensionless Global Error in Synthesis (ERGAS):} Evaluates overall spectral synthesis quality across all bands normalized by band means (lower is better).
    \item \textbf{Correlation Coefficient (CC):} Average Pearson correlation between reconstructed and ground-truth bands (higher is better, max 1.0).
\end{enumerate}
All metrics are evaluated strictly inside the cloud-occluded spatial region defined by binary cloud mask $\mathbf{M} > 0.5$ across all 1,611 held-out test patches.

\begin{table*}[!t]
\centering
\caption{Quantitative Performance Benchmark and Literature Context.}
\label{tab:main_benchmark}
\resizebox{\textwidth}{!}{%
\begin{tabular}{lcccccc}
\toprule
Method & Paradigm / Input Stream & PSNR\,(dB) $\uparrow$ & SSIM $\uparrow$ & SAM\,(deg) $\downarrow$ & ERGAS $\downarrow$ & CC $\uparrow$ \\
\midrule
\multicolumn{7}{l}{\textbf{Direct Experimental Evaluation on 5.0\,m LISS-IV Test Set (1,611 Held-Out Patches):}} \\
Bicubic Spatial Inpainting & Single-Image Interpolation & 6.21 & 0.386 & 10.35 & 166.90 & 0.327 \\
Temporal Baseline (Sentinel-2) & Dry-Season Clear-Sky Prior & 13.11 & 0.656 & 7.48 & 73.99 & 0.397 \\
\textbf{RelieF-CR (Proposed)} & \textbf{Quad-Modal STN Cross-Attention} & \textbf{29.10} & \textbf{0.963} & \textbf{1.49} & \textbf{7.89} & \textbf{0.863} \\
\midrule
\multicolumn{7}{l}{\textit{Literature Reference on 10.0\,m SEN12MS-CR Benchmark (Cited for Architectural Context)$^\dagger$:}} \\
SAR-Opt-cGAN \cite{grohnfeldt2018conditional} & ResNet + cGAN (SAR+Opt) & 24.80 & 0.815 & 5.20 & 18.40 & 0.742 \\
DSen2-CR \cite{meraner2020cloud} & Deep Residual Fusion (SAR+Opt) & 26.35 & 0.864 & 4.15 & 14.20 & 0.781 \\
GLF-CR \cite{xu2022glf} & Global-Local Transformer (SAR+Opt) & 28.40 & 0.908 & 3.42 & 10.50 & 0.812 \\
UnCRtainTS \cite{ebel2023uncrtaints} & Uncertainty Cross-Attention (Multi-Temp+SAR) & 28.95 & 0.924 & 2.85 & 9.80 & 0.835 \\
DiffCR \cite{li2024diffcr} & Conditional Diffusion Prior (SAR+Opt) & 29.20 & 0.931 & 2.40 & 8.90 & 0.841 \\
\bottomrule
\end{tabular}%
}
\vspace{1mm}
\raggedright{\footnotesize $^\dagger$External literature figures are reported on the standard 10\,m SEN12MS-CR benchmark dataset and are provided for architectural context; direct baseline reproduction on high-resolution LISS-IV is subject to sensor-specific adaptation.}
\end{table*}

As detailed in Table~\ref{tab:main_benchmark}:
\begin{itemize}
    \item \textbf{Bicubic Inpainting} fails dramatically (PSNR $6.21\,\text{dB}$, SSIM $0.386$, ERGAS $166.90$) due to the inability of spatial interpolation to reconstruct large opaque cloud gaps.
    \item \textbf{Temporal Substitution} suffers from severe spectral mismatch (PSNR $13.11\,\text{dB}$, SAM $7.48^\circ$) caused by seasonal agricultural phenology and illumination shifts.
    \item \textbf{RelieF-CR} achieves \textbf{29.10\,dB PSNR} and \textbf{0.963 SSIM} on the held-out 5\,m test set, successfully combining structural radar backscatter, clear temporal priors, and topographic relief descriptors.
    \item \textbf{Spectral Angle Mapper (SAM):} Our model achieves \textbf{1.49$^\circ$}, confirming strict radiometric preservation and spectral fidelity across VNIR bands for downstream vegetation modeling.
\end{itemize}

\subsection{Qualitative Analysis Under Severe Cloud Coverage}
In Fig.~\ref{fig:qualitative_grid}, we examine qualitative performance under varying terrain types and cloud densities. Even under severe cloud obstruction ($>75\%$ cloud cover in Fig.~\ref{fig:qualitative_grid}, Row 2), RelieF-CR successfully recovers parcel boundaries, irrigation channels, and forest canopy textures without generating phantom blur or chromatic aberration.
